% Part: history
% Chapter: biographies 
% Section: georg-cantor

\documentclass[../../../include/open-logic-section]{subfiles}

\begin{document}

\olfileid{his}{bio}{can}

\olsection{જ્યૉર્જ કૅન્ટૉર}

\olphoto{cantor-georg}{જ્યૉર્જ કૅન્ટૉર}

જ્યૉર્જ કૅન્ટૉર (ઉચ્ચાર: ગે-ઓર્ગ કાન-ટોર)ના
એક પ્રારંભિક જીવનચરિત્રમાં
દાવો કરાયો હતો કે રશિયાના સેન્ટ પીટર્સબર્ગ તરફ
જતા જહાજમાં તેમનો જન્મ થયો, તેઓ ત્યાં મળી આવ્યા,
અને તેમના માતાપિતા કોણ હતા તે અજાણ્યું હતું.
પરંતુ આ સાચું નથી; તેમનો જન્મ 1845માં સેન્ટ
પીટર્સબર્ગમાં થયો હતો.

કૅન્ટૉરે 1867માં બર્લિન યુનિવર્સિટીમાંથી
ગણિતમાં ડૉક્ટરેટ મેળવી. ગણસિદ્ધાંતમાંના તેમના
કાર્ય માટે તેઓ જાણીતા છે, અને ગણસિદ્ધાંતને
અલગ સંશોધનક્ષેત્ર તરીકે સ્થાપવાનો શ્રેય તેમને
આપવામાં આવે છે. જુદા કદના અનંત ગણો હોય છે
તે સૌપ્રથમ તેમણે સાબિત કર્યું. તેમના સિદ્ધાંતો,
ખાસ કરીને અનંતતાઓનો સિદ્ધાંત, તે સમયના
ગણિતજ્ઞોમાં ઘણી ચર્ચાનું કારણ બન્યા અને તેમનું
કાર્ય વિવાદાસ્પદ રહ્યું.

કૅન્ટૉરની ધાર્મિક માન્યતાઓ અને તેમનું ગાણિતિક
કાર્ય અવિભાજ્ય રીતે જોડાયેલાં હતાં; તેમણે
એવો પણ દાવો કર્યો કે અનંતાતીત સંખ્યાઓનો
સિદ્ધાંત ઈશ્વરે તેમને સીધો જણાવ્યો હતો.
જીવનના પાછળના સમયગાળામાં કૅન્ટૉરને માનસિક
બીમારી થઈ. 1894થી અને પછીનાં વર્ષોમાં વધુ
વાર તેમને હોસ્પિટલમાં દાખલ કરાયા. તેમના
કાર્યની કઠોર ટીકા, જેમાં ગણિતજ્ઞ લિયોપોલ્ડ
ક્રોનેકર સાથેનો સંબંધવિચ્છેદ પણ હતો, તેમને
અવસાદ તરફ અને ગણિતમાં રસ ઘટવા તરફ દોરી ગઈ.
અવસાદના પ્રસંગોમાં કૅન્ટૉર તત્ત્વચિંતન અને
સાહિત્ય તરફ વળતા. તેમણે ફ્રાન્સિસ બેકન
શેક્સપિયરના નાટકોના લેખક હતા એવો સિદ્ધાંત
પણ પ્રકાશિત કર્યો.

કૅન્ટૉરનું અવસાન 6 જાન્યુઆરી 1918ના રોજ
હાલેના એક સેનેટોરિયમમાં થયું.

\begin{reading}
કૅન્ટૉરનાં વિગતવાર જીવનચરિત્રો માટે
\citet{Dauben1990} અને
\citet{Grattan-Guinness1971} જુઓ. તેમના
ક્રાંતિકારી વિચારોનું વર્ણન બીબીસી રેડિયો 4ના
\emph{A Brief History of
  Mathematics} કાર્યક્રમમાં પણ છે
\citep{Sautoy2014}. કૅન્ટૉરના સિદ્ધાંતો
રેપ સ્વરૂપે સાંભળવા હોય તો \citet{Rose2012} જુઓ.
\end{reading}

\end{document}
